\documentclass[a4paper]{article}
\usepackage{amsfonts}
\usepackage{amsmath}
\usepackage{amssymb}
\usepackage{graphicx}     

\newcommand{\dint}{\displaystyle\int}
\newcommand{\llim}{\lim\limits}

\begin{document}
  
\noindent \large Auteur: Oubayy Ahale          \\
\noindent \large Studierichting: Informatica   \\
\noindent \large Oefeningenreeks 1C nr. 46.2   \\

\medskip

\normalsize

$ z^4 = -1 $\\

$ r = |z| = \sqrt{(-1)^2 + 0^2} = \sqrt{1} = 1 $\\

$ \tan(\theta) = \dfrac{0}{-1} = 0             $\\

$ \theta = 0 $\\

We tellen hierbij $\pi$ op want onze $a$ is negatief: \\

$ \Rightarrow \theta = \pi $\\

\medskip

$ \Rightarrow z = 1 \cdot \operatorname{cis}(\pi) =  \operatorname{cis}(\pi) $\\

$w_0 =  \operatorname{cis}\left(\dfrac{\pi}{4}\right) = \cos\left(\dfrac{\pi}{4}\right) + i \cdot \sin\left(\dfrac{\pi}{4}\right) = \dfrac{\sqrt{2}}{2} + \dfrac{\sqrt{2} \cdot i}{2} $\\

$w_1 =  \operatorname{cis}\left( \dfrac{\pi + 2\pi}{4} \right) = \cos\left(\dfrac{3\pi}{4}\right) + i \cdot \sin\left(\dfrac{3\pi}{4}\right) = \dfrac{-\sqrt{2}}{2} + i \cdot \dfrac{\sqrt{2}}{2} $\\

$w_2 =  \operatorname{cis}\left(\dfrac{\pi + 4\pi}{4}\right) =  \operatorname{cis}\left(\dfrac{5\pi}{4}\right) = \cos\left(\dfrac{5\pi}{4}\right) + i \cdot \sin\left(\dfrac{5\pi}{4}\right) = \dfrac{-\sqrt{2}}{2} - \dfrac{\sqrt{2} \cdot i}{2} $\\

$w_3 =  \operatorname{cis}\left(\dfrac{\pi + 6\pi}{4}\right) = \operatorname{cis}\left(\dfrac{7\pi}{4}\right) = \cos\left(\dfrac{7\pi}{4}\right) + i \cdot \sin\left(\dfrac{7\pi}{4}\right) = \dfrac{\sqrt{2}}{2} - \dfrac{\sqrt{2} \cdot i}{2} $\\

\begin{thebibliography}{999}
%\bibitem{a} Referentie
\end{thebibliography}
\end{document} 